\documentclass[10pt, landscape]{article}
\usepackage[scaled=0.92]{helvet}
\usepackage{calc}
\usepackage{multicol}
\usepackage[a4paper,margin=3mm,landscape]{geometry}
\usepackage{amsmath,amsthm,amsfonts,amssymb}
\usepackage{color,graphicx,overpic}
\usepackage{hyperref}
\usepackage{newtxtext} 
\usepackage{enumitem}
\usepackage[table]{xcolor}
\usepackage{mathtools}
\setlist{nosep}
% for including images
\graphicspath{ {../images/} }

\pdfinfo{
  /Title (CS4221.pdf)
  /Creator (TeX)
  /Producer (pdfTeX 1.40.0)
  /Author (Gavin Chiam)
  /Subject (CS4221 Midterm Cheatsheet)
/Keywords (CS4221, nus, cheatsheet, pdf)}

% Turn off header and footer
\pagestyle{empty}

% redefine section commands to use less space
\makeatletter
\renewcommand{\section}{\@startsection{section}{1}{0mm}%
  {-1ex plus -.5ex minus -.2ex}%
  {0.5ex plus .2ex}%x
{\normalfont\large\bfseries}}
\renewcommand{\subsection}{\@startsection{subsection}{2}{0mm}%
  {-1explus -.5ex minus -.2ex}%
  {0.5ex plus .2ex}%
{\normalfont\normalsize\bfseries}}
\renewcommand{\subsubsection}{\@startsection{subsubsection}{3}{0mm}%
  {-1ex plus -.5ex minus -.2ex}%
  {1ex plus .2ex}%
{\normalfont\small\bfseries}}%
\makeatother

\renewcommand{\familydefault}{\sfdefault}
\renewcommand\rmdefault{\sfdefault}
%  makes nested numbering (e.g. 1.1.1, 1.1.2, etc)
\renewcommand{\labelenumii}{\theenumii}
\renewcommand{\theenumii}{\theenumi.\arabic{enumii}.}
\renewcommand\labelitemii{•}
\renewcommand\labelitemiii{•}

\definecolor{mathblue}{cmyk}{1,.72,0,.38}
\everymath\expandafter{\the\everymath \color{mathblue}}

% Don't print section numbers
\setcounter{secnumdepth}{0}

\setlength{\parindent}{0pt}
\setlength{\parskip}{0pt plus 0.5ex}
%% adjust spacing for all itemize/enumerate
\setlength{\leftmargini}{0.5cm}
\setlength{\leftmarginii}{0.5cm}
\setlist[itemize,1]{leftmargin=2mm,labelindent=1mm,labelsep=1mm}
\setlist[itemize,2]{leftmargin=4mm,labelindent=1mm,labelsep=1mm}
\setlist[itemize,3]{leftmargin=4mm,labelindent=1mm,labelsep=1mm}

% adding my commands
\input{../../commands/style-helpers.tex}

% -----------------------------------------------------------------------

\begin{document}
\raggedright
\footnotesize
\begin{multicols*}{4}
  % multicol parameters
  \setlength{\columnseprule}{0.25pt}

  \begin{center}
    \fbox{%
      \parbox{0.8\linewidth}{\centering \textcolor{black}{
          {\Large\textbf{CS4221}}
        \\ \normalsize{AY25/26 SEM 2}}
        \\ {\footnotesize \textcolor{gray}{github/Gavino3o}}
      }%
    }
  \end{center}

  \section{SQL Tuning}

  % Add here after completing later sections

  \section{Functional Dependencies}

  \begin{itemize}
    \item Functinoal dependencies are denoted by $\sigma$
    \item Given: $X \rightarrow Y$.\\If two tuples of $r$ (a table with relation schema $R$) agress on their X-values, then they agree on their Y-values.
    \item Note: SQL query cannot be used in \texttt{CHECK} constraints
    \item Functional dependencies may be enforced when LHS has \texttt{PRIMARY KEY} or \texttt{UNIQUE} constraints.
    \item \textbf{Trivial} FD: $X \rightarrow Y$ where $Y \subseteq X$
    \item \textbf{Non-trivial} FD: $X \rightarrow Y$ where $Y \subsetneq X$
    \item \textbf{Completely non-trivial} FD: $X \rightarrow Y$ where $Y \neq \emptyset$ and $Y \cap X = \emptyset$
    \item Course convention allows: $\emptyset \rightarrow \emptyset, \emptyset \rightarrow Y, X \rightarrow \emptyset$
    \item \textbf{Superkey}: $X$ is a superkey if $X \rightarrow R$ (i.e. $X$ functionally determines all attributes in $R$)
    \item \textbf{Candidate key}: $X$ is a candidate key if $X$ is a superkey and there is no proper subset of $X$ that is a superkey. I.e $X \rightarrow R$ and $\nexists Y \subsetneq X$ such that $Y \rightarrow R$, $X$ is minimal.
    \item \textbf{Logical Entailment}: $F \models f$ if every relation that satisfies $F$ also satisfies $f$. E.g., if $F = \{X \rightarrow Y, Y \rightarrow Z\}$, then $F \models X \rightarrow Z$.
  \end{itemize}

  \subsection{Armstrong's Axioms}
  \begin{itemize}
    \item \textbf{Reflexivity}: If $Y \subseteq X$, then $X \rightarrow Y$ (i.e. trivial FDs are always valid)
    \item \textbf{Augmentation}: If $X \rightarrow Y$, then $XZ \rightarrow YZ$ for any $Z$ (You can add the same attributes to both sides of an FD and it still holds)
    \item \textbf{Transitivity}: If $X \rightarrow Y$ and $Y \rightarrow Z$, then $X \rightarrow Z$ (i.e. You can short circuit the arrows)
    \item \textbf{Union Rule}: If $X \rightarrow Y$ and $X \rightarrow Z$, then $X \rightarrow YZ$
    \item \textbf{Decomposition Rule}: If $X \rightarrow YZ$, then $X \rightarrow Y$ and $X \rightarrow Z$
    \item \textbf{Pseudotransitivity}: If $X \rightarrow Y$ and $YW \rightarrow Z$, then $XW \rightarrow Z$
    \item \textbf{Sound}: If an FD can be derived from Armstrong's Axioms, then it is logically entailed by the set of FDs.
    \item \textbf{Complete}: If an FD is logically entailed by a set of FDs, then it can be derived from Armstrong's Axioms.
  \end{itemize}

  \subsection{Closure of FDs and Attributes}
  \begin{itemize}
    \item The closure of a set of functional dependencies $\sum$ (denoted $\sum^+$) is the set of all FDs that can be logically entailed by $\sum$.
    \item $\sum^+ = \{ X \rightarrow Y | \sum \models X \rightarrow Y \}$
    \item \textbf{Attribute Closure}: The closure of a set of attributes $X$ under a set of FDs $\sum$ (denoted $X^+$) is the set of all attributes that are functionally determined by $X$ under $\sum$.
    \item $X^+ = \{ A | \exists (X \rightarrow \{A\}) \in \sum^+ \}$
    \item Algorithm to compute \textbf{Attribute Closure}:
    \begin{enumerate}
        \item Initialize $X^+ = X$
        \item For each FD $Y \rightarrow Z$ in $\sum$, if $Y \subseteq X^+$, then add $Z$ to $X^+$
        \item Repeat step 2 until no more attributes can be added to $X^+$
    \end{enumerate}
    \item \textbf{Equivalence of FD sets}: Two sets of FDs $\sum$ and $\sum'$ are equivalent if $\sum^+ = \sum'^+$ (i.e. they have the same closure).
    \item To check for equivalence, check if they cover each other:
    \begin{enumerate}
        \item For each FD in $\sum_A$, use FDs in $\sum_B$ to compute the attribute closure of the LHS of the FD. 
        \item If the closure contains the RHS, then the FD is implied by $\sum_B$. 
        \item If all FDs in $\sum_A$ are implied by $\sum_B$, then $\sum_A$ is implied by $\sum_B$. 
        \item Repeat the same process to check if $\sum_B$ is implied by $\sum_A$. If both are implied, then $\sum_A$ and $\sum_B$ are equivalent.
    \end{enumerate}
  \end{itemize}

  \subsection{Minimal Cover}
  \begin{itemize}
    \item A set of FDs $\sum$ is a minimal if:
    \begin{enumerate}
        \item The \textbf{RHS} of every FD in $\sum$ is a single attribute (i.e. FDs are in canonical form)
        \item The \textbf{LHS} of every FD in $\sum$ is minimal (i.e. there are no extraneous attributes in the LHS)
        \item There are \textbf{no redundant} FDs in $\sum$ (i.e. there is no FD in $\sum$ that can be removed without changing the closure of $\sum$)
    \end{enumerate}
    \item Algorithm to compute minimal cover:
    \begin{enumerate}
        \item Decompose FDs so that the RHS of every FD is a single attribute (i.e. put FDs in canonical form)
        \item For each FD, check if there are extraneous attributes in the LHS. If there are, remove them. To check if an attribute $A$ in the LHS of an FD $X \rightarrow Y$ is extraneous, compute the attribute closure of $X - A$ under the set of FDs (including $X \rightarrow Y$). If the closure contains $Y$, then $A$ is extraneous and can be removed from the LHS.
        \item For each FD, check if it is redundant. To check if an FD $X \rightarrow Y$ is redundant, compute the closure of $X$ under the set of FDs without $X \rightarrow Y$. If the closure contains $Y$, then $X \rightarrow Y$ is redundant and can be removed from the set of FDs.
    \end{enumerate}
    \item Minimal cover is not unique. There can be multiple minimal covers for a given set of FDs.
    \item \textbf{Compact:} There is no FD with the same LHS.
    \item \textbf{Canonical Cover}: Compact and minimal cover. 
  \end{itemize}
  
  \subsection{The Chase Algorithm}
  \begin{itemize}
    \item Test whether a give FD $f$ is logically entailed by a set of FDs $\sum$.
    \item Construct a tableau with two rows and $n$ columns, where $n$ is the number of attributes in the relation schema. 
    \item The first row is filled with $\alpha$ and the second row filled with distinct elements.
    \item For the FD $f: X \rightarrow Y$, replace the elements in the second row with $\alpha$ for all attributes in $X$.
    \item Apply the FDs in $\sum$ to the tableau by changing distinct elements into $\alpha$ until no more changes can be made. Note: LHS column $\alpha$.
  \end{itemize}

  \section{Normalization}

  \section{}

\end{multicols*}

\end{document}
